\section{Interleaved determinants: changing the arithmetic without losing positivity}
\label{sec:interleaved-case}

The principal determinant in the companion paper has the same row and
column nodes. Its primitive values grow. Changing the basis of that
matrix cannot help, but changing the node sets changes the determinant
polynomial itself. Interleaved borders are a natural test: retain a
common block carrying the target $\theta$, and add distinct row and
column nodes to create new cancellations at primes. This geometry was
suggested by Litsey's trigamma construction~\cite{Litsey}. The estimates
below concern the beta kernel defined here, not the constants in that
work. They develop record \recordid{P06}.

\subsection{The matrix and its exact analytic benchmark}
For positive integers $a$ define
\[
 S_\ell(a)=\sum_{j=0}^{a-1}\bigl((4j+1)^{-\ell}-(4j+3)^{-\ell}\bigr),
\]
and the rational pencil
\begin{equation}\label{eq:frame-entry}
 \Phi_{ab}(X)=
 \begin{cases}
 \displaystyle\frac{bS_3(b)-aS_3(a)}{6(b^2-a^2)},&a\ne b,\\[5pt]
 \displaystyle X-S_4(a)+\frac{S_3(a)}{12a},&a=b.
 \end{cases}
\end{equation}
Fix $0<c\leq1$, put $m=\lfloor cr/2\rfloor$, $N=r+m$ and $M=r+2m$,
and take the ordered node sets
\[
 \mathcal A=\{1,\ldots,r\}\cup\{r+1,r+3,\ldots,r+2m-1\},\qquad
 \mathcal B=\{1,\ldots,r\}\cup\{r+2,r+4,\ldots,r+2m\}.
\]
Write $P_{r,c}(X)=\det[\Phi_{ab}(X)]_{a\in\mathcal A,b\in\mathcal B}$.
This formula alone specifies every rational coefficient and permits
exact reproduction; no numerical value of $\theta$ is needed.

Here is the analytic input from the companion paper, stated with its
normalization~\cite[Proposition 3.1 and Theorem 3.3]{Main}:
\begin{equation}\label{eq:frame-kernel}
 \Phi_{ab}(\theta)=\int_0^\infty
       \frac{\omega(s)}{(a^2+s^2)(b^2+s^2)}\dd s,\qquad
 \omega(s)=\frac{\pi^2s^2}{96}(1+A(s)-2A(s)^3),
\end{equation}
where $A(s)=\sech(2\pi s)$, and
\begin{equation}\label{eq:frame-comparison}
 L^N\det\left[\frac1{a+b}\right]
 \leq P_{r,c}(\theta)\leq U^N\det\left[\frac1{a+b}\right],
\end{equation}
with $L=(\pi^3/192)/(1+1/(\sqrt2\pi))^2$ and
$U=(\pi^3/192)(1+2/(3\sqrt6))$.
Both constants are positive and independent of all nodes. Positivity
also follows directly from~\eqref{eq:frame-kernel}: expand the
integral determinant as the integral of the product of its two Cauchy
determinants. For ordered nodes the product is positive off coincidences.
The coefficient of $X^r$ is the disjoint-border minor, positive by the
same argument (and equal to one if $m=0$). Thus $P_{r,c}$ has degree
exactly $r$, and $P_{r,c}(\theta)>0$.

The Cauchy product evaluates its leading real decay. Put
\[
 I(a,b)=\frac{(a+b)^2\log(a+b)+(a-b)^2\log|a-b|
                    -2a^2\log a-2b^2\log b}{2},
 \qquad 0^2\log0=0,
\]
and
\begin{equation}\label{eq:frame-energy}
 E(c)=\frac{I(1,1)+2I(1,1+c)+I(1+c,1+c)}4.
 \qquad
 \log P_{r,c}(\theta)=-E(c)r^2+o(r^2).
\end{equation}
To verify the limit, the row and column counting measures divided by
$r$ both converge to
$\rho_c=(\mathbf1_{[0,1]}+\mathbf1_{[0,1+c]})/2$.
In the Cauchy product, the two Vandermonde numerators supply the
$\log|x-y|$ integral and the denominator supplies $\log(x+y)$.
The leftover scale factor is $-N\log r=o(r^2)$.
Truncate the logarithms at distance $\epsilon$ from their singularities.
The bounded grid density bounds the omitted normalized sums by
$O(\epsilon|\log\epsilon|)+o(1)$; then let $\epsilon$ decrease to zero.
Finally two integrations give
\[
 I(a,b)=\int_0^a\int_0^b\log\frac{x+y}{|x-y|}\dd x\dd y.
\]
The $O(N)$ error in~\eqref{eq:frame-comparison} does not change the limit.
Thus there is no unknown leading analytic gain in this specified family.

\subsection{The arithmetic quantity that would suffice}
Let $D_{r,c}$ be the least positive integer clearing the coefficients
of $P_{r,c}$, let $g_{r,c}$ be the gcd of the coefficients of
$D_{r,c}P_{r,c}$, and put $\kappa_{r,c}=D_{r,c}/g_{r,c}$.
A sufficient irrationality criterion for one fixed $c$ is
\begin{equation}\label{eq:frame-target}
 \limsup_{r\to\infty}\frac{\log\kappa_{r,c}}{r^2}<E(c).
\end{equation}
Indeed, rationality $\theta=u/v$ would make
$v^r\kappa_{r,c}P_{r,c}(\theta)$ a positive integer; its logarithm
would tend to minus infinity. Replacing $\kappa_{r,c}$ by $D_{r,c}$
gives a stronger sufficient condition, not an equivalent one.

The following calculation explains both the available estimate and
its limitation. For a prime $p$, the Gauss valuation of a polynomial
is the minimum valuation of its coefficients. It extends to a discrete
valuation on $\Q_p(X)$; reduction of integral rational functions gives
the residue field $\mathbb F_p(X)$. We use elementary pivot elimination
over this valuation ring. This is the setting for the Smith exponents
below; it does not specialize $X$ to $\theta$.

\subsection{Isolating the four negative prime layers}
Suppose $p>3$ and $p^2>4M$. Separate the multiples of $p$ in each tail:
\begin{align*}
 H_\ell(a;p)&=\sum_{\substack{t>0\ {
m odd}\\pt<4a}}
                  \frac{\chi_4(t)}{t^\ell},&
 U_\ell(a;p)&=\sum_{\substack{0<n<4a\ {
m odd}\\p\nmid n}}
                  \frac{\chi_4(n)}{n^\ell},\\
 S_\ell(a)&=\chi_4(p)p^{-\ell}H_\ell(a;p)+U_\ell(a;p).
\end{align*}
All denominators in $H_\ell$ are $p$-units. In a block of $4p$
integers, odd positions differing by $2p$ have the same residue modulo
$p$ and opposite characters. Reflecting a complete block by
$n\mapsto4p-n$ gives, respectively,
\[
 U_3(a+p)\equiv U_3(a),\qquad
 U_3(p-a)\equiv-U_3(a),\qquad U_3(pk)\equiv0\pmod p.
\]
Consequently the matrix obtained by replacing $S_\ell$ by $U_\ell$
in~\eqref{eq:frame-entry} is $p$-integral. When $a\equiv\pm b\ne0$
the numerator cancels the factor of $p$ in $b^2-a^2$; when
$a\equiv b\equiv0$, both factors of $p$ are cancelled by the node
and $U_3$ factors. The inequality $2M<p^2/2$ excludes higher node-factor
valuations. The diagonal follows in the same way.

We therefore have $\Phi=\Psi+R(X)$ with $R$ integral and
\[
 \Psi_{ab}=\frac{\chi_4(p)(bH_3(b;p)-aH_3(a;p))}
                    {6p^3(b^2-a^2)}\quad(a\ne b),\qquad
 \Psi_{aa}=\chi_4(p)\left(-\frac{H_4(a;p)}{p^4}
                              +\frac{H_3(a;p)}{12ap^3}\right).
\]
In particular $p^4\Phi\equiv p^4\Psi\pmod{p^4}$.
Let $\mu_i$ be the Smith exponents of $p^4\Phi$, capped at four,
and $\tau_p=\sum_i(4-\mu_i)$. Successive unit-pivot elimination,
dividing out the minimum power of $p$ at each step, gives
\begin{equation}\label{eq:frame-residual}
 P_{r,c}=p^{-\tau_p}u_p\det T_p,
 \qquad v_G(u_p)=0,\quad T_p\text{ integral}.
\end{equation}
The matrix $T_p$ is what remains after the negative layers have been
removed. Thus $v_G(P_{r,c})=-\tau_p+t_p$, where
$t_p=v_G(\det T_p)\geq0$. Knowing all four negative layers does not
bound $t_p$ and does not determine primitive numerator content.

\subsection{An explicit large-prime calculation}
\begin{proposition}[Four-layer estimate; P06]\label{prop:frame-layers}
For $2M+4<p<4r$, put $h=r-\lfloor p/4\rfloor$. The capped Smith
exponents of $p^4\Phi$ are
\[
 0\ (h\text{ times}),\quad1\ (2m\text{ times}),\quad
 2\ (h\text{ times}),\quad4\ (N-2h-2m\text{ times}).
\]
For $\max(2M+4,4r)<p<4M$, put $H=M-\lfloor p/4\rfloor$.
They are $1$ repeated $H$ times and $4$ repeated $N-H$ times.
Consequently the contribution of these primes to $\sum_p\tau_p\log p$
is $(3+3c^2)r^2+o(r^2)$.
\end{proposition}
\begin{proof}
All node sums and differences are $p$-units in this range. The only
possible odd multiple of $p$ in a tail is $p$ itself, so
$H_3(a;p)=H_4(a;p)=\mathbf1_{4a>p}$. Call these nodes high.
There are $h$ high common nodes and all borders are high in the first
range. Up to the unit $\chi_4(p)$ the singular matrix is $D+pF$,
where $D$ is minus the identity on those $h$ diagonals and zero elsewhere.
The low-low block of $F$ is zero; its low-high entry is
$b/[6(b^2-a^2)]$, a Cauchy entry times a unit column factor.
Every admissible maximal minor is a unit. There are
$l=r-h\geq h+m$ low common nodes.

Eliminate the $h$ unit pivots. On the remaining low and border nodes,
the first layer has block form
$\left(\begin{smallmatrix}0&C_1\\C_2&C_3\end{smallmatrix}\right)$,
where $C_1$ has column rank $m$ and $C_2$ has row rank $m$.
Its rank is $2m$. Left and right null vectors can be supported on the
low common nodes and chosen to annihilate the border Cauchy columns
and rows exactly. On these kernels the second layer is, up to a sign,
\[
 (YF_{\rm low,high\ common})(F_{\rm high\ common,low}Z).
\]
Adjoining the $h$ high-common columns to the $m$ border columns is a
unit full-rank Cauchy matrix, and the analogous assertion holds for
rows. The first factor is injective and the second surjective on an
$h$-dimensional space; their product has rank $h$. Eliminating the
first layer changes this block only in the third layer.
The exact singular matrix $D+pF$ has rank $2(h+m)$: its low-low
block is zero and its two low-high blocks have full rank $h+m$.
The $h+2m+h$ pivots found account for this rank, so there are no
further pivots of valuation less than four.

In the second range no common node is high. The low-high Cauchy
blocks now give rank $H$, all in the first layer; the same argument
leaves every other capped exponent equal to four.
The respective values of $\tau_p$ are $6h+6m$ and $3H$.
With $x=p/r$, the prime number theorem gives the two integrals
\[
 \int_{2+2c}^4 6(1+c/2-x/4)\dd x=3(1-c^2),\qquad
 \int_4^{4+4c}3(1+c-x/4)\dd x=6c^2.
\]
Bounded endpoint strips and the floors contribute $o(r^2)$.
\end{proof}

This estimate cannot by itself prove~\eqref{eq:frame-target}, for any
fixed $0<c\leq1$. The energy integrand is nonnegative and
$\rho_c\leq1$, so
\begin{equation}\label{eq:frame-gap}
 E(c)\leq2(1+c)^2\log2
       \leq4(1+c^2)\log2<3(1+c^2).
\end{equation}
Even omitting every other prime leaves an unfavorable comparison for
the bound obtained by setting $t_p=0$ in~\eqref{eq:frame-residual}.
This proves failure of that bound, not growth of the actual primitive value.
For completeness, using $4N$ at the remaining primes above $\sqrt{4M}$
gives
\[
 \log D_{r,c}\leq(11+12c+7c^2)r^2+o(r^2).
\]
At smaller primes every entry costs $O(\log r/\log p)$ powers of $p$,
so their total is $O(r^{3/2}\log r)=o(r^2)$. This includes 2 and 3.

\subsection{A genuine obstruction for integer multipliers}
For $2M+4<p<4r$ a specific coefficient has the exact valuation
\begin{equation}\label{eq:frame-coefficient}
 v_p([X^{r-h}]P_{r,c})=-4h-3m.
\end{equation}
Indeed, expand this coefficient using minors containing $h$ common
nodes and all borders. A minor containing only $j<h$ high common nodes
has valuation at least $-4j-3(h+m-j)>-4h-3m$.
In the unique minor using all high common nodes, only terms selecting
all $h$ diagonals attain the minimum. The remaining border determinant,
after scaling by $p^{3m}$, is a unit multiple of the Cauchy determinant
$[1/(a+b)]$; its node differences and sums are units. Hence there
is no cancellation of the term of minimum valuation.
Summing~\eqref{eq:frame-coefficient} gives
\[
 \log D_{r,c}\geq(2-c-c^2)r^2+o(r^2).
\]
At $c=1/5$, this lower exponent is $44/25=1.76$, whereas the explicit
formula~\eqref{eq:frame-energy} gives $1.6523<E(1/5)<1.6524$.
The latter interval can be checked using
$\log x=2\sum_{j\geq0}z^{2j+1}/(2j+1)$, $z=(x-1)/(x+1)$,
with geometric tail bounds. Every positive integer coefficient-clearing
multiplier therefore produces growing values at this width.
Division by common numerator content remains possible, so this is
still not a primitive-value obstruction.

\subsection{What remains open, and how to reproduce the construction}
The nine archived examples ($r=8,16,24$ and $c=1/5,1/2,1$) remain
finite evidence. Each used the entire rational polynomial, its least
clearing denominator and the complete coefficient gcd. They do not
bound $g_{r,c}$ for general $r$.
A reader can reproduce such a calculation from~\eqref{eq:frame-entry}
using the block matrix on common and border nodes. If its blocks are
$XI_r+U,V,W,Z$, then $Z$ is invertible by positivity and
\[
 P_{r,c}(X)=\det Z\,\det(XI_r+U-VZ^{-1}W).
\]
Compute the characteristic polynomial over $\Q$, clear all coefficient
denominators, and divide the complete integer coefficient gcd.
This also specifies a practical exact procedure independently of our scripts.

The unresolved quantity is the primitive multiplier $\kappa_{r,c}$.
A new estimate must control the residual factors $t_p$ in
\eqref{eq:frame-residual}, or common content at other primes, strongly
enough to reach~\eqref{eq:frame-target}. The four-layer estimate,
the integer-clearing obstruction and the finite primitive failures
have different scopes; none closes that question.
